\chapter{Topologia da rede e alocação do corte entre usinas}
\label{cap:rede}

A seção~\ref{sec:fig3} testou cinco covariáveis locais ao ponto de conexão e não
achou nenhuma. Restam as covariáveis estruturais, que só existem se houver um
grafo de transmissão montado: distância até a carga, grau do nó, número de
usinas no ponto e outras cinco.

\section{Oito covariáveis de rede na geração fotovoltaica}
\label{sec:fig10}

\begin{ficha}
  \item[Janela]      setembro de 2025 a agosto de 2026
  \item[Amostra]     59 conjuntos fotovoltaicos
  \item[Unidade]     conjunto
  \item[Fonte]       EPE, linhas e subestações em operação; ONS
  \item[Teste]       Spearman sobre 8 covariáveis, com correção de Holm
\end{ficha}

A rede da EPE traz 920 subestações e 1.896 linhas resolvidas em aresta, ou 87\%
do total, com ancoragem mediana de 0,10 km entre a coordenada da usina e o nó
mais próximo. As oito covariáveis foram pré-declaradas antes do teste, e Holm
foi aplicado sobre as oito.

\begin{figure*}[t]
\centering
\includegraphics[width=\linewidth]{../figures/fig10_topologia_rede.pdf}
\caption{\textbf{A distância de rede até a demanda está associada à alocação do
corte que covariáveis locais não explicavam.}
\textbf{a}, Grafo de transmissão em operação, com subestações e linhas da EPE
separadas em 500 kV ou mais e abaixo disso, e as usinas fotovoltaicas da amostra
posicionadas no seu ponto de conexão. A cor codifica a taxa de corte na janela
de referência e o tamanho, a capacidade instalada. As estrelas marcam as nove
regiões metropolitanas usadas como centros de carga. A moldura cobre a região
que contém a amostra; 353 das 1.896 linhas ficam fora dela e não são desenhadas,
e todas as arestas participam do cálculo dos caminhos mínimos.
\textbf{b}, Correlação de Spearman entre a taxa de corte e cada uma das oito
covariáveis de rede pré-declaradas. Em vermelho as que sobrevivem à correção de
Holm.
\textbf{c}, Taxa de corte contra a distância de rede ao centro de carga mais
próximo, por subsistema, com a mediana por quartil de distância.
\textbf{d}, As duas covariáveis sobreviventes medidas dentro de cada subsistema,
para separar efeito de rede de efeito de subsistema.}
\label{fig:rede}
\end{figure*}

Duas covariáveis sobrevivem a Holm, com magnitude próxima e sinais opostos. A
distância de rede à carga dá $\rho = +0{,}38$ e a capacidade da usina dá
$\rho = -0{,}38$, ambas com $p = 0{,}023$ após correção. A soma de tensões no nó
atinge $p < 0{,}05$ bruto e não sobrevive à correção. As outras cinco não
atingem.

\conclusao{Quanto maior a distância em quilômetros de linha até o centro de
carga mais próximo, maior a taxa de corte. A estrutura local do nó não apresenta
associação.}

O resultado é consistente com a seção~\ref{sec:fig9}, em que o Nordeste é
cortado com o corredor operando abaixo do seu percentil habitual. As duas
medidas apontam para a distância até a demanda, e não para a saturação local do
ponto de conexão.

\campo{Não é apenas efeito de subsistema} A distância mediana à carga é próxima
nos dois subsistemas, com 497 km no Nordeste e 520 no Sudeste com Centro-Oeste.
A associação se mantém dentro do Nordeste, com $\rho = +0{,}34$, $p = 0{,}027$ e
$n = 43$. Ela não se mantém dentro do Sudeste com Centro-Oeste, onde dá
$\rho = +0{,}31$, $p = 0{,}25$ e $n = 16$. A capacidade da usina se mantém nos
dois, com $-0{,}41$ e $-0{,}44$, e no Nordeste com $p = 0{,}006$.

\campo{Uma sobrevivente sem hipótese direcional} A capacidade da usina tem sinal
negativo: as maiores são cortadas menos. A direção não havia sido prevista. As
explicações plausíveis, como qualidade de conexão e seleção contratual, são
testáveis e não foram testadas.

\begin{objecao}
A distância de rede não é elétrica: é caminho ponderado por quilômetro de linha.
A impedância está em dado aberto, em 100\% das linhas do pacote
\texttt{linha-transmissao}, com capacidade em MVA em 86,7\% delas. A escolha da
camada geométrica da EPE impediu o uso da distância elétrica nesta seção. Os
centros de carga são nove metrópoles, um proxy geográfico declarado. São oito
hipóteses sobre $n = 59$: Holm foi aplicado, e não há poder para efeitos
pequenos. O resultado é associação, e não causalidade.
\end{objecao}

\campo{Posição na literatura} Maji et al. \cite{maji2025} sustentam que o corte é
fenômeno de nó, e que estimá-lo exige descer ao ponto de conexão. É a hipótese
que esta seção testa e, na geração fotovoltaica, não confirma. O preditor é a
distância até a carga, uma grandeza de caminho, e não de nó. Newbery
\cite{newbery2024} fornece o enquadramento econômico: se a posição na rede
determina a exposição ao corte e o regime de acesso não a precifica, o sinal de
investimento aponta para o lugar errado.

\section{Replicação na geração eólica}
\label{sec:fig11}

\begin{ficha}
  \item[Janela]      novembro de 2022 a julho de 2025, eólica;
                     setembro de 2025 a agosto de 2026, solar
  \item[Amostra]     120 conjuntos eólicos; 59 solares
  \item[Unidade]     conjunto
  \item[Fonte]       ONS, \textit{constrained-off} eólico e fotovoltaico; EPE
  \item[Teste]       as mesmas 8 covariáveis, com Holm por fonte
\end{ficha}

O ONS apura \textit{constrained-off} eólico desde outubro de 2021, em 199
conjuntos, contra 87 do lado solar. A replicação usa as mesmas oito covariáveis,
na mesma ordem, sobre o mesmo grafo, com Holm reaplicado dentro de cada fonte.

\begin{figure*}[t]
\centering
\includegraphics[width=\linewidth]{../figures/fig11_replicacao_eolica.pdf}
\caption{\textbf{O resultado da seção anterior não se replica em eólica, e a
dissociação acompanha a origem declarada do corte de cada fonte.}
\textbf{a}, Correlação de Spearman entre a taxa de corte e as oito covariáveis
de rede pré-declaradas, medidas separadamente em fotovoltaica e eólica. A cor
cheia marca as que sobrevivem à correção de Holm aplicada dentro de cada fonte.
\textbf{b}, Energia cortada decomposta por razão e origem da restrição, conforme
classificação operativa do ONS, para cada fonte.
\textbf{c}, As duas covariáveis-chave medidas apenas no Nordeste, onde as duas
amostras coexistem, para separar tecnologia de geografia.
\textbf{d}, Taxa de corte eólica contra o número de conjuntos no mesmo ponto de
conexão, com a mediana por valor.}
\label{fig:eolica}
\end{figure*}

A distância de rede à carga, que na solar dá $+0{,}38$ e sobrevive a Holm, dá
$-0{,}07$ em 120 conjuntos eólicos. A amostra eólica é o dobro da solar, de modo
que a ausência não é falta de poder.

As covariáveis que sobrevivem em eólica são as que não sobreviveram em solar. O
número de usinas no mesmo nó vai de $+0{,}07$ para $+0{,}47$, e o grau do nó vai
de $-0{,}17$ para $+0{,}35$. As duas sobrevivem a Holm.

A classificação de origem do ONS, independente de todas as covariáveis testadas,
registra corte solar 91\% sistêmico e corte eólico 76\%, com 24\% de origem
local, ou 2,7 vezes a fração do solar. A confiabilidade local responde por 7\%
do corte solar e por 22\% do eólico.

\conclusao{A covariável que prevê o corte de cada fonte acompanha a origem
declarada do corte dessa fonte: grandeza sistêmica onde o corte é sistêmico,
grandeza local onde há corte local.}

\campo{Não é geografia} A amostra eólica é 91\% nordestina, e a solar não é.
Restringindo as duas ao Nordeste, a solar mantém distância à carga em $+0{,}34$
e usinas no nó em $+0{,}05$, enquanto a eólica apresenta $-0{,}19$ e $+0{,}36$.
A dissociação se mantém sob controle de geografia.

\campo{O que muda na seção anterior} A seção~\ref{sec:fig10} enunciou a hipótese
como sobre posição na rede, independentemente da tecnologia. A replicação
restringe o enunciado a fonte cujo corte é majoritariamente sistêmico. O
resultado permanece na geração fotovoltaica, e o que muda é a generalidade.

\campo{Consequência operativa} A exposição solar responde a aproximação da
demanda, por armazenamento ou deslocamento de carga. A exposição eólica responde
a reforço local de escoamento. Uma política única erraria em uma das duas.

\begin{objecao}
As janelas não coincidem: eólica de novembro de 2022 a julho de 2025, solar de
setembro de 2025 a agosto de 2026. Parte da diferença pode ser de período, e o
teste restrito ao Nordeste controla geografia, e não tempo. A eólica se aglomera
mais, com mediana de 4 conjuntos por nó contra 2 no solar, o que facilita
detectar aglomeração. A distância à carga, porém, tem faixa semelhante nas duas
fontes, e não prevê o corte eólico. Os regimes diferem nesta amostra, com 8\% de mediana
em eólica e 26\% em solar, contra 15,9\% e 22,1\% nas frotas completas. Origem é
classificação do ONS, e não inferência causal.
\end{objecao}

\campo{Posição na literatura} Bird, Lew e Milligan \cite{bird2016} organizam a
revisão de onze países pelo eixo entre excesso de oferta sistêmico e restrição
local. A EirGrid \cite{eirgrid} formaliza a distinção no vocabulário, com
\textit{curtailment} para razão sistêmica e \textit{constraint} para local, e
reporta 6,6\% contra 4,7\% no \textit{dispatch-down} eólico de 2025. Não foi
localizado teste publicado de se a covariável preditiva do corte varia com a
composição da causa, e a busca não foi sistemática.

\section{Robustez dos quatro resultados}
\label{sec:fig12}

\begin{ficha}
  \item[Janela]      as amostras das duas seções anteriores
  \item[Amostra]     59 conjuntos solares; 120 eólicos
  \item[Unidade]     conjunto
  \item[Fonte]       ONS, Rede de Operação por número de barra; EPE
  \item[Teste]       bootstrap com 4.000 reamostras; $z$ de Fisher
\end{ficha}

O ONS publica a rede com conectividade por número de barra, impedância em 100\%
das linhas e capacidade em MVA em 87\%. Isso permite reconstruir o grafo por via
independente da camada geométrica usada nas duas seções anteriores. Os três
critérios são o intervalo bootstrap que exclui zero, a estabilidade entre as
metades de uma partição estrutural e a reprodução no grafo alternativo.

\begin{figure*}[t]
\centering
\includegraphics[width=\linewidth]{../figures/fig12_robustez.pdf}
\caption{\textbf{Dos quatro resultados de rede, um passa nos três critérios e o
principal da seção~\ref{sec:fig10} passa em um.}
\textbf{a}, Distribuição bootstrap do coeficiente de Spearman de cada resultado,
com 4.000 reamostras com reposição. A barra marca o intervalo de 95\% e o ponto,
o valor da amostra.
\textbf{b}, Coeficiente calculado separadamente nas duas metades de uma partição
estrutural, entre usinas cuja barra de conexão consta do cadastro de rede de
230 kV ou mais do ONS e as demais, com o valor de $p$ do teste $z$ de Fisher
para a diferença entre metades.
\textbf{c}, O mesmo coeficiente calculado no grafo da EPE e no grafo do ONS
construído por número de barra. Abaixo de cada barra, o tamanho da amostra no
grafo alternativo, ou a marca de que o critério não se aplica.
\textbf{d}, Quantos dos critérios aplicáveis cada resultado passa. O traço marca
o critério não aplicável.}
\label{fig:robustez}
\end{figure*}

\begin{table}[t]
\centering
\tabfont
\begin{tabular}{@{}>{\rag\arraybackslash}p{3.4cm}ccc@{}}
\toprule
\cab{Resultado} & \cab{IC} & \cab{estável} & \cab{outro grafo} \\
\midrule
eólica, grau do nó         & sim & sim & sim \\
eólica, usinas no nó       & sim & não & sim \\
solar, capacidade da usina & sim & sim & --- \\
solar, distância à carga   & sim & não & não \\
\bottomrule
\end{tabular}
\caption{Critérios de robustez aplicados aos quatro resultados que sobreviveram
a Holm.}
\label{tab:robustez}
\notas{O traço indica critério não aplicável: a capacidade da usina não deriva
do grafo e não tem versão alternativa. Para ela valem dois critérios, e não
três.}
\end{table}

Sobreviveram a Holm quatro resultados, e não três. A capacidade da usina em
solar, com $\rho = -0{,}38$ e $p = 0{,}023$ após correção, tem magnitude igual à
da distância à carga e não constava do conjunto avaliado antes.

A distância à carga tem intervalo bootstrap de 0,12 a 0,61, uma largura de 0,49,
cerca do dobro dos demais. Partindo a amostra, o coeficiente cai de 0,62 para
0,11 conforme a metade, com $p = 0{,}025$, e a metade de maior coeficiente
contém todas as 16 usinas do Sudeste. Isso reabre o confundimento regional que o
painel \textbf{d} da seção~\ref{sec:fig10} afastava.

A capacidade da usina passa nos dois critérios aplicáveis, com intervalo de
$-0{,}60$ a $-0{,}13$ e partição não rejeitada, de $-0{,}18$ contra $-0{,}48$ e
$p = 0{,}21$.

\conclusao{O grau do nó na geração eólica passa nos três critérios. A distância
à carga na geração fotovoltaica passa em um dos três, e é o resultado de rede
menos sustentado.}

Dois critérios não são três, e a diferença entre metades ainda reduz o
coeficiente da capacidade da usina em dois terços. Passar nos critérios
aplicáveis é não ser rejeitado, e não é ser confirmado.

\campo{Amostra, e não método} No grafo do ONS a distância dá $-0{,}14$ em
quilômetros e $-0{,}06$ em reatância. Rodando o grafo da EPE nas mesmas 28
usinas, ele também não apresenta associação. A diferença é de amostra, o que não
recupera o resultado e identifica em que subconjunto ele ocorre.

A distância elétrica foi calculada, e a afirmação da seção~\ref{sec:fig10} de
que ela não estaria disponível é incorreta. A correlação entre a distância
elétrica e a distância em quilômetros é de 0,53, de modo que a substituição
constitui teste independente. O resultado não muda.

\campo{Os critérios não dependiam de fonte nova} Os três critérios usam apenas
dados já disponíveis na seção~\ref{sec:fig10}, e poderiam ter sido aplicados
ali. Foram executados a partir da troca de grafo.

\begin{objecao}
A partição não é aleatória, por construção: ela testa dependência a
característica estrutural, e não variabilidade amostral genérica. O grafo do ONS
é exato onde cobre, e exclui conexões abaixo de 230 kV; para a amostra solar
isso remove mais da metade dos conjuntos e elimina um subsistema inteiro. Três
critérios não esgotam robustez, e faltam validação temporal, análise de
influência ponto a ponto e sensibilidade à definição de centro de carga. Os
limiares são convenção declarada, e o critério do intervalo foi corrigido de
``é positivo'' para ``exclui zero''. Como o grafo é variável derivada de
sensoriamento, ele propaga erro para o intervalo \cite{proctor2023}, e nenhuma
imputação foi aplicada.
\end{objecao}
